% 仙后座（综述）
% license CCBYSA3
% type Wiki

本文根据 CC-BY-SA 协议转载翻译自维基百科\href{https://en.wikipedia.org/wiki/Cassiopeia_(constellation)}{相关文章}。

仙后座（Cassiopeia，ⓘ）是位于北天的一個星座與星群，其名稱來自希臘神話中虛榮的王后卡西奧佩婭（Cassiopeia），她是安德洛墨達（Andromeda）的母親，曾誇耀自己擁有無與倫比的美貌。仙后座是公元 2 世紀希臘天文學家托勒密列出的 48 個星座之一，至今仍是現代 88 個星座之一。由於其由五顆亮星組成的獨特「W」形狀，仙后座極易辨認。

仙后座位於北天，在北緯 34° 以上的地區全年可見；在（亞）熱帶地區，9 月至 11 月初觀測條件最佳；而在南半球低緯度、低於南緯 25° 的熱帶地區，則可於特定季節在北方低空看到。

視星等為 2.2 的仙后座 α 星，即舍達爾（Schedar），是仙后座中最亮的恆星。該星座擁有一些已知最為明亮的恆星，包括黃色特超巨星仙后座 ρ 星（Rho Cassiopeiae）與 V509 Cassiopeiae，以及白色特超巨星仙后座 6 星。1572 年，第谷·布拉赫（Tycho Brahe）所觀測到的超新星就在仙后座中明亮爆發$^{[5]}$。仙后座 A 是一個超新星殘骸，在頻率高於 1 GHz 時，是全天最亮的銀河系外射電源。已有 14 個恆星系統被發現擁有系外行星，其中之一——HD 219134——被認為可能擁有六顆行星。一段富集的銀河帶穿過仙后座，其中包含多個疏散星團、年輕而明亮的銀河盤恆星以及星雲。IC 10 是一個不規則星系，是目前已知距離最近的星暴星系，也是本星系群中唯一的此類星系。
\subsection{神話}
\begin{figure}[ht]
\centering
\includegraphics[width=6cm]{./figures/99da126f8a82dd17.png}
\caption{《乌拉尼娅之镜》中描绘的坐在宝座上的仙后座（Cassiopeia）} \label{fig_Cape_1}
\end{figure}
该星座以埃塞俄比亚王后卡西奥佩娅（Cassiopeia）命名。卡西奥佩娅是埃塞俄比亚国王刻甫斯（Cepheus）的妻子$^{[5]}$，也是公主安德洛墨达（Andromeda）的母亲。刻甫斯与卡西奥佩娅以及安德洛墨达一同被安置在星空中，相互毗邻。她之所以被放置在天空中，作为惩罚，是因为她夸口说自己的女儿安德洛墨达比涅瑞伊得斯（Nereids，海中女神）更加美丽，或者另一种说法是，她自称比海中仙女本身更加美丽，从而激怒了波塞冬$^{[6]}$。她被迫坐在宝座上绕北天极旋转，一半时间必须紧紧抓住宝座以免跌落；而波塞冬则裁定安德洛墨达应被绑在岩石上，作为海怪刻托斯（Cetus）的祭品。随后，安德洛墨达被英雄珀尔修斯所拯救，并最终与其成婚$^{[7][8]}$。

在历史上，卡西奥佩娅星座的形象有多种不同的描绘方式。在波斯，她被苏菲（al-Sufi）描绘为一位女王，右手持有一根带新月的权杖，头戴王冠，身旁还有一只双峰骆驼；在法国，她则被描绘为坐在大理石宝座上，左手持棕榈叶，右手提着自己的长袍。这一形象出自奥古斯丁·罗耶（Augustin Royer）于1679年出版的星图集$^{[7]}$。

在中国天文学中，构成仙后座的恒星分布于三个区域之中：紫微垣（紫微垣，Zǐ Wēi Yuán）、北方玄武（北方玄武，Běi Fāng Xuán Wǔ）以及西方白虎（西方白虎，Xī Fāng Bái Hǔ）。

中国天文学家在现代所称的仙后座区域内辨认出多种星象组合。κ、η 与 μ 仙后座星组成了名为“王桥”的星官；当它们与 α 与 β 仙后座星共同观测时，则构成了名为“王良”的大车。御者手中的鞭子由 γ 仙后座星表示，该星有时被称为“次”，即汉语中“鞭”的意思$^{[7]}$。

在印度神话中，仙后座（Cassiopeia）与神话人物沙尔米什塔（Sharmishtha）相关联——她是伟大的阿修罗（Daitya）国王弗里什帕尔瓦（Vrishparva）的女儿，也是德瓦雅尼（Devayani，对应安德洛墨达）的朋友。

在威尔士神话中，Llys Dôn（字面意思为“多恩的宫廷”）是该星座的传统威尔士名称。多恩的至少三位子女同样具有天文学上的对应：Caer Gwydion（“圭迪恩的堡垒”）是银河的传统威尔士名称，而 Caer Arianrhod（“阿丽安罗德的堡垒”）则对应北冕座（Corona Borealis）$^{[9]}$。

在17世纪，仙后座中的恒星曾被用来描绘多位《圣经》人物，其中包括所罗门之母拔示巴（Bathsheba）、《旧约》中的女先知底波拉（Deborah），以及耶稣的追随者抹大拉的玛利亚（Mary Magdalene）$^{[7]}$。

在一些阿拉伯星图集中，仙后座的恒星还被描绘为一个名为“染色之手”（Tinted Hand）的形象。据不同说法，这只手要么象征一只被指甲花染成红色的女子之手，要么象征穆罕默德之女法蒂玛（Fatima）沾满鲜血的手。该“手”由 α、β、γ、δ、ε 与 η 仙后座星组成；其“手臂”则由 α、γ、δ、ε、η 与 ν 英仙座星组成$^{[7]}$。

另一个包含仙后座恒星的阿拉伯星座是“骆驼”。其头部由仙女座的 λ、κ、ι 与 φ 星构成；驼峰是 β 仙后座星；身体由其余仙后座恒星组成；而双腿则由英仙座与仙女座中的恒星构成$^{[7]}$。

在其他文化中，人们还从这一星群中看到了手或麋鹿鹿角的形象$^{[10]}$。例如在萨米人（Sámi）的传统中，仙后座的 “W” 形状被视为麋鹿的鹿角；西伯利亚的楚科奇人（Chukchi）同样将其中五颗主星视作五只驯鹿雄鹿$^{[7]}$。

马绍尔群岛的居民将仙后座视为一幅巨大鼠海豚星座的一部分：仙后座的主星构成其尾部，仙女座与三角座（Triangulum）构成其身体，而白羊座则构成其头部$^{[7]}$。在夏威夷，α、β 与 γ 仙后座星各自有名称：α 仙后座星被称为 Poloahilani，β 仙后座星被称为 Polula，γ 仙后座星被称为 Mulehu。普卡普卡（Pukapuka）群岛的人们则将仙后座视为一个独立的星座，称为 Na Taki-tolu-a-Mataliki $^{[11]}$。
\subsection{特征}
\begin{figure}[ht]
\centering
\includegraphics[width=6cm]{./figures/8d913297d824a2aa.png}
\caption{夜空中的仙后座} \label{fig_Cape_2}
\end{figure}
仙后座曾发生过一次超新星事件，即仙后座A（Cassiopeia A），SN~1572。

仙后座覆盖面积为 $598.4$ 平方度，占全天面积的 $1.451\%$，在 $88$ 个星座中按面积排名第 $25$ 位$^{[12]}$。它北侧和西侧与仙王座（Cepheus）相邻，南侧和西侧与仙女座（Andromeda）相邻，东南侧与英仙座（Perseus）相邻，东侧与鹿豹座（Camelopardalis）相邻，并在西侧与蝎虎座（Lacerta）共享一小段边界。

该星座的三字母官方缩写由国际天文学联合会（International Astronomical Union）于 1922 年采用，为 ``Cas''$^{[13]}$。官方星座边界由比利时天文学家欧仁·德尔波特（Eugène Delporte）于 $1930$ 年制定$^{[b]}$，以一个由 $30$ 条边组成的多边形加以界定。在赤道坐标系中，这些边界的赤经范围为 $00^{\mathrm h}\,27^{\mathrm m}\,03^{\mathrm s}$ 至 $23^{\mathrm h}\,41^{\mathrm m}\,06^{\mathrm s}$，赤纬范围为 $+77.69^\circ$ 至 $+46.68^\circ$,$^{[3]}$。由于其位于北天球，该星座对南纬 $12^\circ$ 以北的观测者而言均可见$^{[12][c]}$。在高纬北方天空中，它属于拱极星座（即在夜空中永不落下），在不列颠群岛、加拿大以及美国北部地区均可全年观测到$^{[15]}$。
\subsubsection{特征}
\textbf{恒星}
\begin{figure}[ht]
\centering
\includegraphics[width=6cm]{./figures/88f670490ee393b8.png}
\caption{从北半球观测时，肉眼所见的仙后座星座} \label{fig_Cape_3}
\end{figure}
德国制图学家约翰·拜耳（Johann Bayer）使用希腊字母 Alpha 至 Omega，随后又使用 A 和 B，来标记该星座中最显著的 26 颗恒星。后来发现 Upsilon 实际上由两颗恒星组成，并由约翰·弗兰斯蒂德（John Flamsteed）标记为 Upsilon$^{[1]}$ 和 Upsilon$^{[2]}$。B Cassiopeiae 实际上就是著名的超新星——第谷超新星（Tycho's Supernova）$^{[16]}$。在该星座边界之内，共有 157 颗视星等小于或等于 6.5 的恒星$^{[d][12]}$。

“W”星群  
仙后座中最亮的五颗恒星——Alpha、Beta、Gamma、Delta 与 Epsilon Cassiopeiae——构成了具有代表性的 W 形星群$^{[15]}$。这五颗恒星均可用肉眼清晰观测，其中三颗具有明显的变光性，第四颗则被怀疑为低振幅变星。在北半球的春季与夏季夜晚，当该星群位于北极星下方时，其形状呈现为 W；而在北半球冬季，或从南半球纬度观测时，该星群位于北极星“上方”（即更接近天顶），其 W 形看起来则是倒置的。

Alpha Cassiopeiae，传统名称为 Schedar（源自阿拉伯语 Al Sadr，意为“胸部”），常被误认为是一个由四颗恒星组成的系统，但实际上它是一颗主星，伴随三颗在物理上距离遥远的视伴星。主星占据主导地位，是一颗橙色巨星，视星等为 2.2，距离地球$228\pm2$光年$^{[18]}$。其光度约为太阳的 771 倍，在其约 1 至 2 亿年的生命周期中，核心氢燃料耗尽后发生膨胀并冷却；在此前的大部分时间里，它是一颗蓝白色 B 型主序星$^{[19]}$。一颗视星等 8.9 的黄色矮星伴星（B）与其相距甚远，而另外两颗伴星（C 与 D）距离更近，视星等分别为 13 和 14$^{[20]}$。

Beta Cassiopeiae，又称 Caph（意为“手”），是一颗白色恒星，视星等为 2.3，距离地球 $54.7\pm0.3$ 光年$^{[18]}$。其年龄约为 12 亿年，核心氢已耗尽，正开始偏离主序阶段并发生膨胀和冷却。其质量约为太阳的 1.9 倍，光度约为太阳的 21.3 倍。Caph 的自转速度约为其临界速度的 92\%，完成一次自转仅需 1.12 天，这使得恒星呈现为扁球体形状，其赤道半径比极半径大约 24\%$^{[21]}$。它是一颗盾牌座$\delta$ 型变星，振幅较小，周期约为 2.5 小时$^{[22]}$。

Gamma Cassiopeiae，现称为 Tiansi，是 Gamma Cassiopeiae 型变星的原型，这类变星由于高速自转，会将物质抛射形成一个可变的盘状结构。Gamma Cassiopeiae 的最暗视星等为 3.0，最亮可达 1.6，但通常接近 2.2，其亮度变化具有不可预测的减弱与增强。它是一个分光双星系统，轨道周期为 203.59 天，其伴星的计算质量约与太阳相当。然而至今尚未获得该伴星的直接观测证据，因此有人推测它可能是一颗白矮星或其他简并天体$^{[23]}$。该系统距离地球 $550\pm10$ 光年。

Delta Cassiopeiae，也称 Ruchbah 或 Rukbat，意为“膝盖”，可能是一颗 Algol 型食双星，最大亮度约为 2.7 等。有报道称其存在振幅小于 0.1 等、周期约为 $2$ 年 $1$ 个月的食现象$^{[24]}$，但这一结果从未得到确认。它距离地球$99.4\pm0.4$光年$^{[18]}$。Delta Cassiopeiae 曾被让·皮卡尔（Jean Picard）于 1669 年用于测定地球表面上一度弧长，从而推算地球半径$^{[25]}$。

Epsilon Cassiopeiae，又称 Segin，视星等为 3.3，距离地球$410\pm20$光年$^{[18]}$。它是一颗炽热的蓝白色恒星，光谱型为 B3 III，表面温度约为$15{,}680\,\mathrm{K}$。其质量约为太阳的 6.5 倍，直径约为太阳的 4.2 倍，属于 Be 星这一类——高速自转并抛射出环状或壳状物质的恒星$^{[26]}$。

\textbf{较暗恒星}
\begin{figure}[ht]
\centering
\includegraphics[width=6cm]{./figures/f7c9e38a40b34540.png}
\caption{仙后座 $\kappa$ 星及其弓形激波。斯皮策红外望远镜图像（NASA/JPL-Caltech）} \label{fig_Cape_4}
\end{figure}
仙后座中接下来最亮的七颗恒星也全部被确认或被怀疑为变星，其中包括拜耳未赋予希腊字母的 50 仙后座星，它是一颗振幅极小的疑似变星。$\zeta$仙后座星（Fulu$^{[27]}$）是一颗疑似缓慢脉动的 B 型恒星。$\kappa$仙后座星（Cexing）是一颗光谱型为 BC0.7Ia 的蓝色超巨星，其光度约为太阳的 $302{,}000$ 倍，直径为太阳的 33 倍$^{[28]}$。它是一颗高速逃逸星，相对于其邻近恒星以约每小时 2.5 百万英里（每秒 $1{,}100$ 千米）的速度运动$^{[29]}$。其磁场与粒子风在其前方约 4 光年处形成了一个可见的弓形激波，与弥散且通常不可见的星际气体与尘埃发生碰撞。该弓形激波的尺度极其巨大：长度约 12 光年，宽度约 1.8 光年$^{[30]}$。$\theta$仙后座星，名为 Marfak，是一颗疑似变星，其亮度变化小于十分之一个星等。$\iota$仙后座星是一套距地球 142 光年的三星系统，主星是一颗视星等 4.5 的白色恒星，同时也是$\alpha_2$猎犬座型变星；次星是一颗视星等 6.9 的黄色恒星；第三星的视星等为 8.4。主星与次星彼此接近，而主星与第三星则相距甚远。$\omicron$ 仙后座星也是一套三星系统，其主星同样是一颗$\gamma$仙后座型变星。

$\eta$仙后座星（Achird）是仙后座中距离地球最近的恒星，距离约 19.13 光年$^{[31]}$。它是一套双星系统，由一颗与太阳相似、视星等 3.45 的 G 型矮星，以及一颗视星等 7.51 的 K 型矮星组成$^{[32]}$。它在天球上的位置介于$\alpha$与$\gamma$仙后座星之间，是距离地球最近的 G 型／K 型恒星之一。

$\sigma$仙后座星是一套距地球 $1{,}500$ 光年的双星系统，其主星呈绿色、视星等 5.0，次星呈蓝色、视星等 7.3。$\psi$ 仙后座星是一套距地球 193 光年的三星系统，主星是一颗视星等 4.7 的橙色巨星，次星为一对彼此接近、合成视星等约为 $9.0$ 的恒星$^{[24]}$。

$\rho$仙后座星是一颗半规则脉动变星型黄色特超巨星，其光度约为 $500{,}000\,L_{\odot}$，是银河系中最明亮的恒星之一$^{[33]}$。其最暗视星等为 6.2，最亮视星等为 4.1，周期约为 320 天。其直径约为太阳的 450 倍，质量为太阳的 17 倍，初始质量约为太阳的 45 倍。$\rho$ 仙后座星距地球约 $10{,}000$ 光年。仙后座还包括 V509 仙后座星，这是第二颗极为罕见的黄色特超巨星，其光度约为太阳的 $400{,}000$ 倍、质量为太阳的 14 倍$^{[33]}$，以及更高温的白色特超巨星 6 仙后座星。此外，该星座还包含红色超巨星 PZ 仙后座星，它是已知体积最大的恒星之一，估计半径为 $(1{,}190\text{--}1{,}940)\,R_{\odot}$，同时也是一颗半规则变星$^{[34]}$。其光度约为太阳的 $240{,}000$ 至 $270{,}000$ 倍，距离地球约 $9{,}160$ 光年$^{[35]}$。

AO 仙后座星是一套双星系统，由一颗 O8 型主序星和一颗 O9.2 型亮巨星组成，其质量分别约为太阳的 20.30--57.75 倍和 14.8--31.73 倍$^{[36]}$。这两颗大质量恒星彼此距离极近，因而相互引力作用使其形状被拉伸成近似卵形$^{[37]}$。

第谷·布拉赫的新星曾在仙后座中可见，而恒星 Tycho G 被怀疑是为该次爆炸提供触发物质的供体星。
\subsubsection{深空天体}
\begin{figure}[ht]
\centering
\includegraphics[width=6cm]{./figures/b85fc1bca2832008.png}
\caption{行星状星云 IC 289 是一团被恒星核心遗迹向外推入太空的电离气体云。} \label{fig_Cape_5}
\end{figure}
银河系中一段十分丰富的区域贯穿仙后座，从英仙座方向延伸至天鹅座，包含大量疏散星团、年轻而明亮的银河盘恒星以及各类星云。

心脏星云（Heart Nebula）和灵魂星云（Soul Nebula）是彼此相邻的两处发射星云，距离地球约 $7{,}500$ 光年。

仙后座中包含两个梅西耶天体，即 M52（NGC~7654）和 M103（NGC~581），二者均为疏散星团。M52 曾被描述为一个“肾形”的星团，约包含 100 颗恒星，距离地球约 $4{,}600$ 光年$^{[38]}$。其最显著的成员是一颗位于星团边缘、呈橙色、视星等为 8.0 的恒星。M103 的成员数量远少于 M52，仅约 25 颗恒星，同时其距离也更远，约为 $8{,}000$ 至 $9{,}500$ 光年$^{[39]}$。其最显著的成员实际上是一对距离更近、视向叠加的双星系统，由一颗 $7$ 等主星和一颗 10 等伴星组成$^{[24]}$。
\begin{figure}[ht]
\centering
\includegraphics[width=6cm]{./figures/fcc2a9e6beeaf1d2.png}
\caption{仙后座中的四个星团：IC~166、NGC~654、NGC~663 和 NGC~659} \label{fig_Cape_6}
\end{figure}
仙后座中另外一些显著的疏散星团包括 NGC~457 和 NGC~663，二者均约包含 80 颗恒星。NGC~457 的结构较为疏松，其最亮成员为仙后座$\phi$星（Phi Cassiopeiae），这是一颗呈白色、视星等为 5.0 的超巨星。然而，$\phi$仙后座是否真正属于该疏散星团目前尚不确定$^{[40]}$。NGC~457 中的恒星呈链状分布，距离地球约 $10{,}000$ 光年。NGC~663 的距离则更近，约为 $8{,}200$ 光年，同时其视直径更大，约为$0.25^\circ$,$^{[24]}$。

仙后座中存在两个超新星遗迹。第一个编号为 3C~10，亦称第谷超新星遗迹，是被称为“第谷之星”的超新星爆发之后的残留物。该超新星于 1572 年由第谷·布拉赫观测到，如今在射电波段中表现为一个明亮的天体$^{[24]}$。在由仙后座五颗主要恒星构成的 “W” 形星群内部，位于仙后座 A（Cas~A）。它是一次大约 300 年前发生的超新星爆发的遗迹（以地球当前观测到的时间计；其实际距离约为 $10{,}000$ 光年）$^{[41]}$，并且是太阳系之外可观测到的最强射电源之一。该天体可能曾于 1680 年被约翰·弗拉姆斯蒂德记录为一颗暗弱恒星。它也是钱德拉 X 射线天文台在 1990 年代末期获取的首批图像的研究对象之一。由该恒星抛射出的物质外壳以每秒约 $4{,}000$ 千米（$2{,}500$ 英里）的速度膨胀，其平均温度约为 $30{,}000\,\mathrm{K}$,$^{[41]}$。
\begin{figure}[ht]
\centering
\includegraphics[width=6cm]{./figures/8b22206cd9e7fd56.png}
\caption{仙后座，显示了双星团 $\chi$ 和 $h$ 英仙座，以及星团 NGC~654、NGC~663、M~103、NGC~457、NGC~225、NGC~7788、NGC~7790、NGC~7789 和 M~52} \label{fig_Cape_7}
\end{figure}
NGC~457 是仙后座中的另一处疏散星团，也被称为 E.T. 星团、猫头鹰星团以及 Caldwell~13。该星团于 1787 年由威廉·赫歇尔发现。其总星等为 6.4，距离地球约 $10{,}000$ 光年，位于银河系的英仙臂之中。然而，其最显著的成员——双星 $\phi$~Cassiopeiae——要近得多，距离仅约 $1{,}000$ 至 $4{,}000$ 光年。NGC~457 较为富集，属于 Shapley~e 级、Trumpler~I~3~r 型星团，其中心较为集中，并与周围恒星场分离。该星团包含超过 100 颗恒星，其亮度差异很大。$^{[42]}$

仙后座中还包含本星系群的两个成员。NGC~185 是一颗 9.2 等的 E0 型椭圆星系，距离地球约 200 万光年。稍暗且更远的 NGC~147 是一颗 9.3 等的椭圆星系，同样为 E0 型，距离地球约 230 万光年。尽管它们在天空中并不位于仙女座星系之内，但这两个矮星系都在引力上束缚于体量大得多的仙女座星系。$^{[43]}$

IC~10 是一颗不规则星系，是目前已知距离最近的星暴星系，也是本星系群中唯一的此类星系。$^{[44]}$

仙后座还包含距离本星系群最近的星系群之一——IC~342/Maffei 星系群的一部分。Maffei~1 和 Maffei~2 位于心脏星云与灵魂星云以南不远处。由于处在回避带区域，这两者尽管距离均在 $1{,}000$ 万光年以内，却显得异常昏暗（其中 Maffei~2 的亮度低于大多数业余望远镜的观测范围）。$^{[45]}$
\subsubsection{流星雨}  
十二月 $\phi$~仙后座流星雨是一场近年来发现的、发生在十二月初的流星雨，其辐射点位于仙后座。$\phi$~仙后座流星速度极慢，进入大气层时的速度约为每秒 16.7 千米。该流星雨的母体被认为是比拉彗星，但其母体的具体身份曾在相当长一段时间内并不明确。$^{[46]}$
\subsection{名称衍生}  
USS~Cassiopeia（AK-75）是一艘以该星座命名的美国海军陨石坑级货船。

在《宝可梦 朱／紫》中，反派组织“天星队”（Team~Star）被划分为五个小队，其名称均源自仙后座中最亮的恒星：Segin 小队、Schedar 小队、Ruchbah 小队、Navi 小队以及 Caph 小队。该组织的领导者使用“Cassiopeia”这一化名。

在《二之国：白色圣灰的女王》中，倒数第二位主要反派、亦即“白色女巫”，其名字正是卡西欧佩亚女王（Queen~Cassiopeia）。

Cassiopeia 也是伦敦乐队 Bears~in~Trees 的一首歌曲名称。尽管歌词主要指涉古希腊女性卡西欧佩亚，但该专辑封面展示了仙后座星座。$^{[47]}$

Cassiopeia 还是《英雄联盟》中的一名英雄角色。她的美貌与虚荣心映射了希腊神话中的同名人物。

Casiopea 是一个成立于 1976 年的日本爵士融合乐队的名称。该名称由吉他手的母亲选定，唯一的理由是便于乐队成员记住。
\subsection{参见}  
\begin{itemize}
\item 仙后座的恒星形成区
\end{itemize}  
\subsection{参考文献}  
\subsubsection{说明性注释} \\
a.$\gamma$~Cas 为变星，并且在某些时期亮度会超过 $\alpha$。\\
b.Delporte 曾向国际天文学联合会提议将星座边界标准化，该组织同意了这一提议，并授予他主导角色。$^{[14]}$\\  
c.尽管在南纬 $12^\circ$ 至 $43^\circ$ 之间的观测者技术上可以看到星座的部分区域升起至地平线上方，但距离地平线仅数度以内的恒星在实际上几乎无法观测。$^{[12]}$\\  
d.视星等为 6.5 的天体属于在郊区—乡村过渡夜空条件下肉眼可见的最暗天体之一。$^{[17]}$
